\chapter{Architecture: CNNs}
\label{chap:cnn}
We have already introduced convolution's two core assumptions: useful patterns are local and that the same detector can be reused.
So, we will now extend that mental model into efficient vision systems, detection, temporal video and pretrained visual representations.

\section{Detector Perspective}
\label{sec:cnn-local-detectors}
Do not think of convolution as tensor-index arithmetic, think of it as a detector we have tuned.
Imagine each convolutional channel as a detector to certain patterns. 
For example, looking at the ground with many stencils, trying to find animals in grass.
You will drag your grasshopper stencil across the ground in a grid, and when it meets an grasshopper you will have a pattern match (a high output for that grid square), but when it is over the grass there will be no matches.
Once you have traced the grasshopper grid, get out your spider stencil and map out the grid of spiders.
The grid squares are patches and the width of the grid squares are the stride.
Read one filter as a small detector swept over the image:
\[
    \underbrace{\text{feature response at one position}}_{\text{one output value}}
    =
    \sum
    \underbrace{\text{local pixels/features}}_{\text{small receptive patch}}
    \times
    \underbrace{\theta_{\text{shared filter}}}_{\text{same detector everywhere}}.
\]

\begin{figure}[htbp]
    \centering
    \begin{tikzpicture}[x=.67cm,y=.67cm]
    \foreach \x in {0,...,5}{\foreach \y in {0,...,5}{\draw[rulegray] (\x,\y) rectangle ++(1,1);}}
    \fill[palerblue] (1,2) rectangle (4,5);
    \draw[deepblue,very thick] (1,2) rectangle (4,5);
    \node[font=\small,align=center] at (3,6.6) {Image / feature map};
    \node[font=\scriptsize,deepblue] at (2.5,3.5) {local patch};
    \draw[arrblue] (6.4,3.5)--(8.1,3.5);
    \node[flow,minimum width=2.2cm] at (9.5,3.5) {same learned\\filter};
    \draw[arrblue] (10.9,3.5)--(12.5,3.5);
    \node[yesbox] at (13.8,3.5) {response};
    \end{tikzpicture}
    \caption{Convolution is easiest to read as one reusable local detector evaluated at many positions.}
\end{figure}

\paragraph{Channels are Functions}
RGB begins with three maps.
After a learned layer, the channels are better viewed as different feature maps: edge evidence, texture evidence, colour transitions, or more abstract learned patterns.
Later layers mix these channels into progressively richer evidence.


\begin{calloutmemoryhook}
A CNN reuses local pattern detectors across space.
\end{calloutmemoryhook}

\section{Receptive field, resolution and depth}
\label{sec:cnn-receptive-field}
A shallow unit sees a small patch; deeper units combine evidence from many earlier patches and therefore depend on a larger part of the input.
Depth builds a hierarchy from local evidence toward larger-context evidence.

\begin{figure}[htbp]
\centering
\begin{tikzpicture}[node distance=7mm and 9mm]
\node[flow] (p) {Pixels\\small patch};
\node[flow,right=of p] (e) {Edges /\\textures};
\node[flow,right=of e] (s) {Shapes /\\parts};
\node[yesbox,right=of s] (o) {Object-level\\evidence};
\draw[arr] (p)--(e);
\draw[arr] (e)--(s);
\draw[arr] (s)--(o);
\end{tikzpicture}
\caption{Stacking local operations increases effective context; the labels are a useful hierarchy intuition rather than fixed neuron semantics.}
\end{figure}

Stride and pooling reduce spatial resolution.
The practical trade-off is
\[
\underbrace{\text{high resolution}}_{\text{fine localisation, more memory/compute}}
\quad\longleftrightarrow\quad
\underbrace{\text{low resolution}}_{\text{cheaper, coarser spatial detail}}.
\]
Detection and segmentation usually care more about preserving location than image-level classification.

\section{Efficient convolution: separate spatial filtering from channel mixing}
\label{sec:cnn-depthwise-separable}
A standard convolution performs two jobs together: inspect local spatial structure and mix information across channels.
A \textbf{depthwise-separable convolution} factorises them:
\[
\boxed{
\text{spatial filtering per channel}
\ \longrightarrow\
\text{$1\times1$ channel mixing}
}.
\]
For a $3\times3$ layer with $64$ input and $128$ output channels, the rough per-position weight interaction count changes from
\[
\underbrace{3\times3\times64\times128}_{73{,}728\ \text{standard}}
\quad\longrightarrow\quad
\underbrace{3\times3\times64+64\times128}_{8{,}768\ \text{factorised}}.
\]
The reduction comes from a restriction: spatial and cross-channel reasoning are no longer represented by one unrestricted operation.
MobileNet-style architectures exploit this trade-off for resource-constrained inference.

\begin{calloutbias}
If spatial filtering and channel mixing can be approximately factorised, depthwise-separable convolution saves substantial compute; if their interaction must be richly coupled, the restriction can cost accuracy.
\end{calloutbias}

\section{Classification, detection and cheap classical stages}
\label{sec:cnn-classification-detection}
Classification maps an image to an image-level label.
Detection must additionally return locations, so its output is a set of boxes/classes/confidences rather than one class.

\begin{figure}[htbp]
\centering
\begin{tikzpicture}[node distance=8mm and 14mm]
\node[flow] (im1) {Image};
\node[flow,right=of im1] (cls) {Classifier};
\node[yesbox,right=of cls] (lab) {``dog''};
\node[flow,below=14mm of im1] (im2) {Image};
\node[flow,right=of im2] (det) {Detector};
\node[yesbox,right=of det] (boxes) {dog box\\ball box};
\draw[arr] (im1)--(cls);\draw[arr] (cls)--(lab);
\draw[arr] (im2)--(det);\draw[arr] (det)--(boxes);
\end{tikzpicture}
\caption{Classification asks what is present; detection asks what is present and where.}
\end{figure}

A learned model is not automatically the cheapest valid solution.
For a strongly coloured marker, HSV thresholding plus contour geometry can provide a fast first-stage detector:
\[
\text{RGB frame}
\rightarrow
\text{HSV}
\rightarrow
\text{colour threshold}
\rightarrow
\text{contours}
\rightarrow
\text{candidate location}.
\]
A production cascade can therefore begin with deterministic image structure and use a neural detector only for harder appearance variation.

\begin{calloutpractical}
Use the cheapest hypothesis class that captures the structure you actually need; a learned model should earn its extra flexibility.
\end{calloutpractical}

\section{Video: use time when single frames are ambiguous}
\label{sec:cnn-temporal}
Motion blur, occlusion or lighting can make one frame unreliable even when neighbouring frames are clear.
A temporal stage uses nearby observations to stabilise the decision:
\[
\underbrace{\hat y_t}_{\text{current decision}}
=
f(
\underbrace{x_{t-k:t}}_{\text{recent frames}}
).
\]
The useful assumption is that the physical scene evolves continuously enough for neighbouring frames to provide evidence about the present.

\begin{calloutbias}
If nearby frames describe one continuously evolving scene, temporal consistency can suppress transient visual noise; abrupt scene cuts or genuinely fast changes weaken that assumption.
\end{calloutbias}

\section{Pretrained visual representations: CLIP-like versus DINO-like}
\label{sec:cnn-pretrained-visual-representations}
The most important difference is the pretraining signal rather than the brand name.
A CLIP-like system aligns an image encoder with a text encoder:
\[
\text{image}
\rightarrow
z_{\text{image}}
\quad\text{close to}\quad
z_{\text{text}}
\leftarrow
\text{matching caption}.
\]
Its representation is therefore pressured toward linguistically describable semantics.
A DINO-like self-supervised visual encoder learns reusable image structure from transformed views without requiring paired captions.
Both can become backbones for a small downstream classifier; \cref{sec:adaptation-zero-shot-learning-no-task-specific-labelled-examples} gives the zero-shot image--text connection.

\begin{center}
\small
\begin{tabularx}{.94\linewidth}{@{}L{.23\linewidth}YY@{}}
\toprule
\textbf{Representation} & \textbf{Training pressure} & \textbf{Useful consequence} \\
\midrule
CLIP-like & Match paired image/text semantics & Language-aligned similarity and zero-shot prompts \\
DINO-like & Self-supervised consistency across image views & General visual features without captions \\
\bottomrule
\end{tabularx}
\end{center}

\section{How vision systems evolved}
\label{sec:cnn-evolution}
The useful evolution is an accumulation of constraints and reusable representation learning rather than a replacement ladder.

\begin{center}
\small
\begin{tabularx}{.97\linewidth}{@{}L{.18\linewidth}L{.22\linewidth}YY@{}}
\toprule
\textbf{Step} & \textbf{What changed} & \textbf{Problem targeted} & \textbf{New trade-off} \\
\midrule
Classical CV & Hand-written colour/geometry rules & Extremely cheap structured tasks & Brittle to appearance variation \\
CNN & Learn local shared filters & Hand-written feature limits & Training/data/compute cost \\
Residual/deeper CNN & Easier optimisation of depth & Need larger effective hierarchy & More compute \\
Depthwise conv. & Factorise spatial/channel work & Edge/mobile serving cost & More constrained layer \\
Pretrained encoder & Reuse broad learned representation & Limited labelled target data & Inherit pretraining bias \\
Multimodal/self-supervised & Broaden training signal & Reusable semantics without task labels & Larger pretraining systems \\
\bottomrule
\end{tabularx}
\end{center}

\begin{calloutdegreeofexperiment}
Vary input resolution, kernel/receptive-field design, depth, width, standard versus factorised convolution, downsampling schedule, pretrained backbone, temporal context and deployment precision.
\end{calloutdegreeofexperiment}
